% Generated draft: requires mathematical and source review.
\documentclass[11pt]{article}
\usepackage[margin=25mm]{geometry}
\usepackage{fontspec}
\setmainfont{Latin Modern Roman}
\usepackage{amsmath,amssymb,mathtools,unicode-math}
\setmathfont{Latin Modern Math}
\usepackage{enumitem,needspace,xurl,hyperref}
\setlength{\parindent}{0pt}
\setlength{\parskip}{5pt}
\setlength{\emergencystretch}{3em}
\newcommand{\N}{\mathbb{N}}
\newcommand{\Z}{\mathbb{Z}}
\newcommand{\Q}{\mathbb{Q}}
\newcommand{\R}{\mathbb{R}}
\newcommand{\C}{\mathbb{C}}
\newcommand{\F}{\mathbb{F}}
\newcommand{\A}{\mathbb{A}}
\newcommand{\PP}{\mathbb{P}}
\newcommand{\E}{\mathbb{E}}
\newcommand{\eps}{\varepsilon}
\newcommand{\dd}{\,\mathrm{d}}
\begin{document}
\section*{UnsolvedMath ID: 30005221}
\textbf{Source problem code:} catalogue entry
\subsection{Definitions and mathematical statement}
Let \(p\) be a prime, let \(G\) be a finite group, and let \(P\) be a Sylow \(p\)-subgroup of \(G\). Write \(\operatorname{Irr}(G)\) for the irreducible complex characters of \(G\). For \(\chi\in\operatorname{Irr}(G)\), put
\[
\Q(\chi)=\Q\bigl(\chi(g):g\in G\bigr),\qquad \chi_P=\operatorname{Res}^{G}_{P}\chi.
\]
Under the cyclotomic interpretation of the supplied notation, let \(\Q_n=\Q(\zeta_n)\), where \(\zeta_n\) is a primitive \(n\)-th root of unity, and let the conductor of \(\chi\) mean the least positive integer \(n\) such that \(\Q(\chi)\subseteq\Q_n\). Write this conductor as \(p^a m\), with \(a\geq0\), \(m\geq1\), and \(p\nmid m\). Define
\[
\Q_p(\chi_P)=\Q\bigl(\zeta_p,\chi(x):x\in P\bigr),
\qquad
\chi(1)_p=p^{v_p(\chi(1))},
\]
where \(v_p(d)\) is the largest nonnegative integer \(b\) such that \(p^b\mid d\).

\textbf{Question, with an unresolved boundary qualification.} For every such \(p,G,P,\chi\) with \(a\geq1\), does the following inequality hold, without any assumption that \(G\) is \(p\)-solvable?
\[
\bigl[\Q_{p^a}:\Q_p(\chi_P)\bigr]\leq\chi(1)_p.
\]
Here the bracket denotes the degree of a field extension. With the cyclotomic interpretation above and \(a\geq1\), the fields have the inclusion needed to form this degree. A finite group is \(p\)-solvable if it has a finite normal series whose factors each have either \(p\)-power order or order relatively prime to \(p\).

The supplied statement does not explicitly require \(a\geq1\). When \(a=0\), its upper field is \(\Q_1=\Q\), whereas its lower field contains \(\Q_p\); for odd \(p\), the displayed expression therefore does not define an extension degree. The intended treatment of this boundary case is missing. The positive-exponent formulation above must consequently be reviewed against the intended full question. The cyclotomic notation and conductor convention are also interpretations requiring confirmation from the report.

The original excerpt separately asks whether ``the theorem above'' extends to non-\(p\)-solvable groups, but does not supply that theorem. Its full extension question cannot be stated precisely from this record. The excerpt also contains an incomplete discussion of the linear constituents of an unspecified character \(\psi\), grouped by their orders; it does not provide enough context to turn that discussion into an additional assertion or hypothesis.
\subsection{Short English statement}
The question asks whether the degree of a cyclotomic field over the field generated by the restricted character values and a primitive \(p\)-th root of unity is at most the \(p\)-part of the character degree, for every finite group. The supplied formulation requires clarification when the conductor is prime to \(p\), and its separate reference to an earlier theorem is incomplete. This is a statement-only draft requiring human mathematical review.
\subsection{Sources}
\begin{itemize}
\item Supplied local registry record: permanent UnsolvedMath ID \(30005221\), source problem code \texttt{catalogue entry}, title \emph{Extending Character-Field Bounds Beyond p-Solvable Groups}, in the Oberwolfach Reports Open Problems collection. The record gives proposed year 2022 and names no individual proposer.
\item Supplied report citation: \emph{Character Theory and Categorification (2023)}, Oberwolfach Reports, DOI 10.4171/owr/2022/39, pp. 2267--2333. \url{https://doi.org/10.4171/owr/2022/39}. The supplied title year is retained alongside the report identifier 39/2022.
\item Supplied literature source: Fayers--Kleshchev--Morotti, spin RoCK blocks. \url{https://doi.org/10.1112/jlms.12852}. The record does not explain this reference's connection to the character-field inequality.
\item Supplied literature source: Oberwolfach Report 39/2022. \url{https://ems.press/content/serial-article-files/46974}.
\item Literature assessment dated exactly \texttt{2026-08-22}: the supplied curation labels the problem partially solved and states that the Isaacs--Navarro theorem settles the stated \(p\)-solvable case and establishes substantial arithmetic restrictions, while the non-\(p\)-solvable extension and proposed general degree inequality remain unresolved. This is supplied dated curation, not an independently verified assessment. No separate citation or complete statement of the Isaacs--Navarro theorem is supplied.
\item The registry labels its statement \texttt{corrected\_verified} and reports that the official Oberwolfach report was checked, the question isolated from surrounding exposition, and its notation repaired. These are supplied source-verification claims; no independent source visit or mathematical verification is asserted here.
\end{itemize}
\end{document}
